\documentclass[11pt,a4paper]{article}
\usepackage{fontspec,amsmath,amssymb,amsthm,microtype}
\setmainfont{lmroman10-regular.otf}[BoldFont=lmroman10-bold.otf,ItalicFont=lmroman10-italic.otf,BoldItalicFont=lmroman10-bolditalic.otf]
\usepackage{xcolor}
\usepackage[most]{tcolorbox}
\usepackage{mathrsfs}
\providecommand{\lt}{<}
\providecommand{\gt}{>}
\setmonofont{DejaVuSansMono.ttf}[BoldFont=DejaVuSansMono-Bold.ttf,ItalicFont=DejaVuSansMono-Oblique.ttf,BoldItalicFont=DejaVuSansMono-BoldOblique.ttf,Scale=0.86]
\newfontfamily\leanfallback{STIXTwoMath-Regular.otf}[Scale=0.86]
\definecolor{leanrule}{gray}{0.62}
\definecolor{leanback}{gray}{0.962}
\definecolor{leantitle}{gray}{0.22}
\definecolor{leanmuted}{gray}{0.38}
\newtcolorbox{leanbox}[1]{enhanced,breakable,colback=leanback,colframe=leanrule,
  boxrule=0.4pt,arc=0.6mm,left=2.4mm,right=2.4mm,top=1.8mm,bottom=1.8mm,
  before skip=0.8em,after skip=1em,fonttitle=\small\bfseries,coltitle=leantitle,
  colbacktitle=leanback,titlerule=0.3pt,title={#1}}
\newlength{\leancharwidth}
\newenvironment{leancode}{\par\smallskip\begingroup\ttfamily\footnotesize\raggedright
  \settowidth{\leancharwidth}{M}\setlength{\parskip}{0pt}\setlength{\parindent}{0pt}
  \frenchspacing}{\par\endgroup\smallskip}
\newcommand{\leanline}[2]{\hangindent=\dimexpr#1\leancharwidth+4\leancharwidth\relax
  \hangafter=1\noindent\hspace*{#1\leancharwidth}#2\par}
\newcommand{\leansource}[1]{\par{\footnotesize\color{leanmuted}#1}\par}
\newcommand{\leanlabel}[1]{\par\smallskip\noindent{\small\bfseries #1}\par}
\usepackage[margin=25mm]{geometry}
\usepackage[unicode,hidelinks]{hyperref}
\hypersetup{pdfauthor={LangevinLean Project}}
\setlength{\emergencystretch}{3em}
\setlength{\parskip}{0.4em}
\setlength{\parindent}{0pt}
\allowdisplaybreaks[2]
\raggedbottom
\title{\texorpdfstring{The Langevin update and cancellation under shared noise}{The Langevin update and cancellation under shared noise}}
\author{LangevinLean Project}
\date{}
\begin{document}\maketitle
Draft mathematical exposition; these notes may contain mistakes.


For a natural number \(d\), let \(E_d=\mathbb R^d\) with the Euclidean inner product \(\langle x,y\rangle=\sum_{i=0}^{d-1}x_iy_i\) and norm \(\|x\|=\langle x,x\rangle^{1/2}\); it is a real Hilbert space, and for \(d=0\) it is the zero space. For a potential \(U:E_d\to\mathbb R\) with gradient \(\nabla U\), a step size \(h\in\mathbb R\) and a noise vector \(z\in E_d\), one step of the unadjusted Langevin algorithm (ULA) maps \(x\in E_d\) to


\[T_h^z(x)=x-h\nabla U(x)+\sqrt{2h}\,z,\]


with the convention \(\sqrt{2h}=0\) for \(h<0\). If \(h>0\) and \(z\) is a standard Gaussian vector, \(T_h^z(x)\) has the Gaussian law with mean \(x-h\nabla U(x)\) and covariance \(2hI_d\); this is the Euler–Maruyama step for the Langevin equation \(dX_t=-\nabla U(X_t)\,dt+\sqrt2\,dB_t\).




When the same noise vector is used at two states, it cancels from the difference of the two updates, which then equals the difference of two deterministic gradient steps. Combined with the previous section, this gives the shared-noise estimate \(\|T_h^z(x)-T_h^z(y)\|\le\sqrt{1-2hm+h^2L^2}\,\|x-y\|\) for a differentiable \(m\)-strongly convex potential with \(L\)-Lipschitz gradient. The estimate holds for every vector \(z\); no probabilistic assumption on the noise is involved.


\section{\texorpdfstring{The Euclidean state space}{The Euclidean state space}}


\textbf{Definition.} For \(d\in\mathbb N\), \(E_d\) is the space \(\mathbb R^d\) of vectors \(x=(x_0,\dots,x_{d-1})\), with the operations defined coordinatewise, the inner product \(\langle x,y\rangle=\sum_{i=0}^{d-1}x_iy_i\) and the norm \(\|x\|=\bigl(\sum_{i=0}^{d-1}x_i^2\bigr)^{1/2}\).




Being finite-dimensional, \(E_d\) is complete, hence a real Hilbert space; all statements of the previous sections that require a Hilbert space apply to it. For \(d=0\) the index set is empty and \(E_0=\{0\}\).

\begin{leanbox}{Lean: \texttt{E}}
\leansource{\texttt{E} \quad UnadjustedLangevin/\allowbreak{}ULAStep.lean, lines 14–14}
\begin{leancode}
\leanline{0}{abbrev\ E\ (d\ :\ ℕ)\ :=\ EuclideanSpace\ ℝ\ (Fin\ d)}
\end{leancode}
\leanlabel{How the Lean proof follows the argument}
\texttt{abbrev E (d : \ensuremath{\mathbb N}) := Euclidean\allowbreak{}Space \ensuremath{\mathbb R} (Fin d)}. Here \texttt{Fin d} is the set \(\{0,\dots,d-1\}\) and \texttt{Euclidean\allowbreak{}Space \ensuremath{\mathbb R} \ensuremath{\iota}} is Mathlib's space of real vectors indexed by a finite set \(\iota\) with the Euclidean (\(\ell^2\)) inner product and norm. Declaring \texttt{E} as an \texttt{abbrev} makes it reducible, so that the inner-product, completeness and measurability structures of \texttt{Euclidean\allowbreak{}Space} are used for \texttt{E d} without further declarations.
\end{leanbox}

\section{\texorpdfstring{One unadjusted Langevin step}{One unadjusted Langevin step}}


\textbf{Definition.} For \(U:E_d\to\mathbb R\), \(h\in\mathbb R\) and \(x,z\in E_d\), the ULA update is


\[T_h^z(x)=x-h\nabla U(x)+\sqrt{2h}\,z,\]


where \(\sqrt{2h}=0\) if \(h<0\).




The update consists of the gradient step \(x-h\nabla U(x)\) and the noise increment \(\sqrt{2h}\,z\). The definition is a deterministic function of \((x,z)\); probabilistic statements arise only when \(x\) and \(z\) are replaced by random vectors. No sign condition on \(h\) is imposed here; the estimates below assume \(h\ge0\).

\begin{leanbox}{Lean: \texttt{ulaUpdate}}
\leansource{\texttt{ulaUpdate} \quad UnadjustedLangevin/\allowbreak{}ULAStep.lean, lines 17–18}
\begin{leancode}
\leanline{0}{def\ ulaUpdate\ \{d\ :\ ℕ\}\ (U\ :\ E\ d\ →\ ℝ)\ (h\ :\ ℝ)\ (x\ z\ :\ E\ d)\ :\ E\ d\ :=}
\leanline{2}{x\ -\ h\ •\ gradient\ U\ x\ +\ Real.sqrt\ (2\ *\ h)\ •\ z}
\end{leancode}
\leanlabel{How the Lean proof follows the argument}
\texttt{ula\allowbreak{}Update U h x z} is defined as \texttt{x -\allowbreak{} h \ensuremath{\bullet} gradient U x + Real.\allowbreak{}sqrt (2 * h) \ensuremath{\bullet} z}; \texttt{\ensuremath{\bullet}} is scalar multiplication and \texttt{Real.\allowbreak{}sqrt} is \(0\) on negative arguments. \texttt{gradient U x} is the Riesz representative of the derivative of \(U\) at \(x\); Mathlib sets it to \(0\) at points where \(U\) is not differentiable, but every theorem below assumes that \(U\) is differentiable, so this convention never intervenes. The definition is outside the namespace \texttt{ULA}.
\end{leanbox}

\section{\texorpdfstring{Shared noise cancels}{Shared noise cancels}}


\textbf{Claim.} For every \(U:E_d\to\mathbb R\), \(h\in\mathbb R\) and \(x,y,z\in E_d\),


\[T_h^z(x)-T_h^z(y)=\bigl(x-h\nabla U(x)\bigr)-\bigl(y-h\nabla U(y)\bigr).\]


\textbf{Proof.} With \(a=x-h\nabla U(x)\), \(b=y-h\nabla U(y)\) and \(c=\sqrt{2h}\,z\), the two updates are \(a+c\) and \(b+c\), and \((a+c)-(b+c)=a-b\) in any vector space. \(\square\)




The cancellation requires the same vector \(z\) in both updates. With two different noise vectors \(z,z'\), the additional term \(\sqrt{2h}(z-z')\) would remain.

\begin{leanbox}{Lean: \texttt{ulaUpdate\_\allowbreak{}sub}}
\leansource{\texttt{ulaUpdate\_\allowbreak{}sub} \quad UnadjustedLangevin/\allowbreak{}ULAStep.lean, lines 21–24}
\begin{leancode}
\leanline{0}{lemma\ ulaUpdate\_sub\ \{d\ :\ ℕ\}\ (U\ :\ E\ d\ →\ ℝ)\ (h\ :\ ℝ)\ (x\ y\ z\ :\ E\ d)\ :}
\leanline{4}{ulaUpdate\ U\ h\ x\ z\ -\ ulaUpdate\ U\ h\ y\ z\ =}
\leanline{6}{(x\ -\ h\ •\ gradient\ U\ x)\ -\ (y\ -\ h\ •\ gradient\ U\ y)}
\end{leancode}
\leanlabel{How the Lean proof follows the argument}
\texttt{ula\allowbreak{}Update\_\allowbreak{}sub} is proved by unfolding \texttt{ula\allowbreak{}Update} and applying the group identity \texttt{add\_\allowbreak{}sub\_\allowbreak{}add\_\allowbreak{}right\_\allowbreak{}eq\_\allowbreak{}sub}, \((a+c)-(b+c)=a-b\).
\end{leanbox}

\section{\texorpdfstring{The shared-noise estimate from strong monotonicity}{The shared-noise estimate from strong monotonicity}}


\textbf{Claim.} Let \(U:E_d\to\mathbb R\) be differentiable at every point and \(m\)-strongly convex, where \(m\in\mathbb R\), and let \(\nabla U\) be \(L\)-Lipschitz with \(L\ge0\). Let \(h\ge0\). Then for all \(x,y,z\in E_d\),


\[\|T_h^z(x)-T_h^z(y)\|\le\sqrt{1-2hm+h^2L^2}\;\|x-y\|.\]


\textbf{Proof.} By the previous step, the left-hand side equals \(\|(x-h\nabla U(x))-(y-h\nabla U(y))\|\). Apply the square-root Euler-step bound of the previous section with \(g_x=\nabla U(x)\) and \(g_y=\nabla U(y)\). Its two hypotheses hold: \(m\|x-y\|^2\le\langle\nabla U(x)-\nabla U(y),x-y\rangle\) is the strong monotonicity of the gradient, proved for differentiable \(m\)-strongly convex functions on the Hilbert space \(E_d\); and \(\|\nabla U(x)-\nabla U(y)\|\le L\|x-y\|\) is the Lipschitz assumption. \(\square\)




The factor is less than one exactly when \(h(hL^2-2m)<0\), that is, when \(h>0\) and \(hL^2<2m\). The sharp contraction of the next section improves it: if \(0<m\le L\), then \((1-mh)^2\le1-2hm+h^2L^2\) and \((1-Lh)^2\le1-2hm+h^2L^2\), because \(m^2\le L^2\) and \(-2hL\le-2hm\), so the sharp factor \(\max\{|1-mh|,|1-Lh|\}\) is never larger.

\begin{leanbox}{Lean: \texttt{ula\_\allowbreak{}shared\_\allowbreak{}noise\_\allowbreak{}contract}}
\leansource{\texttt{ula\_\allowbreak{}shared\_\allowbreak{}noise\_\allowbreak{}contract} \quad UnadjustedLangevin/\allowbreak{}ULAStep.lean, lines 37–48}
\begin{leancode}
\leanline{0}{theorem\ ula\_shared\_noise\_contract\ \{d\ :\ ℕ\}\ (U\ :\ E\ d\ →\ ℝ)\ (m\ h\ :\ ℝ)\ (L\ :\ ℝ≥0)}
\leanline{4}{(hdiff\ :\ Differentiable\ ℝ\ U)}
\leanline{4}{(hsc\ :\ StrongConvexOn\ Set.univ\ m\ U)}
\leanline{4}{(hsm\ :\ LipschitzWith\ L\ (gradient\ U))}
\leanline{4}{(hh\ :\ 0\ ≤\ h)}
\leanline{4}{(x\ y\ z\ :\ E\ d)\ :}
\leanline{4}{‖ulaUpdate\ U\ h\ x\ z\ -\ ulaUpdate\ U\ h\ y\ z‖\ ≤}
\leanline{6}{Real.sqrt\ (1\ -\ 2\ *\ h\ *\ m\ +\ h\ \textasciicircum{}\ 2\ *\ (L\ :\ ℝ)\ \textasciicircum{}\ 2)\ *\ ‖x\ -\ y‖}
\end{leancode}
\leanlabel{How the Lean proof follows the argument}
\texttt{ula\_\allowbreak{}shared\_\allowbreak{}noise\_\allowbreak{}contract} takes \texttt{hdiff : Differentiable \ensuremath{\mathbb R} U}, \texttt{hsc : Strong\allowbreak{}Convex\allowbreak{}On Set.\allowbreak{}univ m U}, \texttt{hsm : Lipschitz\allowbreak{}With L (gradient U)} and \texttt{hh : 0 \ensuremath{\le} h}. After \texttt{rw [ula\allowbreak{}Update\_\allowbreak{}sub]} (the noise cancellation of the previous step), the goal is closed by \texttt{ULA.\allowbreak{}gradient\_\allowbreak{}step\_\allowbreak{}contract} applied to \texttt{gradient U x} and \texttt{gradient U y}. Its monotonicity hypothesis is \texttt{ULA.\allowbreak{}strong\_\allowbreak{}gradient\_\allowbreak{}monotone hdiff hsc x y}; its Lipschitz hypothesis is \texttt{hsm.\allowbreak{}dist\_\allowbreak{}le\_\allowbreak{}mul x y}, with distances rewritten as norms of differences by \texttt{dist\_\allowbreak{}eq\_\allowbreak{}norm}.
\end{leanbox}

\end{document}
